\begin{abstract}
본 논문은 멀티-툴 도구 선택에서 실패하는 것은 저랭크 steering 자체가 아니라 \emph{정지 K-style steering의 배치 방식}라는 주장을 제시한다. MetaTool Subtask4에서 정지 K는 $-4.57$pp로 붕괴하지만, 같은 온톨로지 기저 위에서 K를 초기 레이어에만 제한하고 Q를 전역에 두는 \emph{layer-adaptive K+Q}는 $+2.08$pp를 회복한다. $\tau^2$-bench에서도 이 설계는 retail/telecom/airline에서 각각 $+5.98$/$+26.76$/$+3.84$pp를 기록하며, retail을 action 수로 분해하면 short-horizon에서는 layer-adaptive가, medium/long-horizon에서는 signed Q-only가 우세해 두 연산자가 같은 operator family 안의 regime-dependent 양끝임이 드러난다. Llama-3.1-8B-Instruct에서도 telecom under-focused regime에서는 $+11.62$pp가 재현되지만, 최적 signed-Q 부호는 Qwen과 반대이고 positive $\beta$는 대규모 empty output collapse를 유발한다. 이론적으로 우리는 정지 K가 방출 히스토리를 인코딩할 수 없음을 보이고, K-측과 Q-측 개입이 per-step에서는 동등하지만 KV-cache 재사용 아래 multi-step에서는 갈라짐을 증명한다. 요점은 단순하다. 멀티-툴에서 중요한 것은 ``더 강한 정지 steering''이 아니라 \emph{어느 깊이에서 어떤 쪽을 조작할 것인가}다.
\end{abstract}
